\chapter{ĐẶC TẢ BẢNG, HÌNH VÀ NỘI DUNG CẦN CẬP NHẬT}
\label{app:pending-assets}

\section{Tệp hình Draw.io dự kiến}
Mỗi hình minh họa có một tên tệp PDF cố định. LaTeX tự hiển thị placeholder nếu chưa có file, vì vậy CI vẫn build được. Khi tạo bằng Draw.io, xuất cả SVG và PDF vào thư mục \texttt{thesis/figures/}.

\begin{table}[H]
\caption{Danh sách hình minh họa cần tạo bằng Draw.io}
\label{tab:drawio-assets}
\begin{tabularx}{\textwidth}{p{1.0cm}>{\raggedright\arraybackslash\fontsize{10.5}{13}\selectfont}p{6.6cm}Y}
\toprule
\textbf{Ch.} & \textbf{Tệp PDF} & \textbf{Mục đích}\\
\midrule
1 & \path{chapter1_research_landscape.pdf} & Bản đồ tài liệu, phương pháp và khoảng trống\\
2 & \path{chapter2_conceptual_foundation.pdf} & Liên kết RSMA--STAR--RIS--TD3\\
3 & \path{chapter3_system_model.pdf} & Sơ đồ BS, STAR--RIS, user truyền/phản xạ và đường kênh\\
4 & \path{chapter4_td3_pipeline.pdf} & Pipeline ScenarioBank, train, validation, test và audit\\
5 & \path{chapter5_quality_latency_tradeoff.pdf} & Sơ đồ Pareto định tính giữa quality, QoS và latency\\
\bottomrule
\end{tabularx}
\tablesource{Tác giả xây dựng (2026)}
\end{table}

\section{Các bảng kết quả nên thay bằng raw bundle khi chốt bản nộp}
Bản luận văn hiện sử dụng các giá trị đã được xác nhận trong reviewer summary và không phát sinh giá trị trung gian. Trước bản nộp cuối, nên tự động sinh các bảng sau từ raw CSV:
\begin{enumerate}[label=\arabic*)]
  \item mean, standard deviation và 95\% CI của sum-rate theo method và $N$;
  \item QoS fraction, all-QoS và violation theo $N$;
  \item mean/median/p95 latency;
  \item paired differences, effect size và adjusted $p$-value;
  \item ablation theo $N$;
  \item convergence và best checkpoint step theo seed.
\end{enumerate}

\section{Quy tắc không tạo số liệu}
Không được nội suy giá trị cho $N=32,64,96$ từ hai endpoint $N=16$ và $N=128$. Không được chuyển range thành thứ tự theo $N$ nếu raw table chưa xác nhận. Không được báo DDPG/PPO là đã chạy đầy đủ nếu thiếu manifest. Mọi bảng mới phải chứa nguồn ``Tác giả tổng hợp từ raw evidence'' và lưu script sinh bảng.

\section{Thông tin bìa cần xác nhận}
Các macro ở đầu \texttt{main.tex} hiện sử dụng thông tin từ đề cương: tác giả Nguyễn Duy Thanh, mã học viên 2480101831, lớp KM2402, ngành Khoa học máy tính, mã ngành 8480101, hai người hướng dẫn và năm 2026. Trước khi in, cần xác nhận cách gọi chính thức ``Luận văn thạc sĩ'' hay ``Đề án tốt nghiệp thạc sĩ'', tháng nộp, logo chuẩn và cách viết tên người hướng dẫn.
